\documentclass[11pt,a4paper]{article}
\usepackage[utf8]{inputenc}
\usepackage[T1]{fontenc}
\usepackage{times}
\usepackage{amsmath,amssymb}
\usepackage{booktabs}
\usepackage{multirow}
\usepackage{graphicx}
\usepackage{xcolor}
\usepackage[margin=2.5cm]{geometry}
\usepackage{natbib}
\usepackage{hyperref}
\usepackage{tablefootnote}

\newcommand{\lzstar}{$l_{z}^{*}$}
\newcommand{\lzm}{$l_{z}^{m}$}
\newcommand{\thetafix}{$\theta$-fixed}

\title{IRT-FIT: Detecting Benchmark Contamination in LLMs\\from Binary Response Patterns Alone}

\author{Anonymous}

\date{}

\begin{document}
\maketitle

% ============================================================
\begin{abstract}
Benchmark contamination threatens the reliability of LLM evaluation, yet every existing detection method requires privileged access beyond binary (correct/incorrect) responses---training data overlap, token probabilities, model weights, or specially crafted prompts.
We introduce \textbf{IRT-FIT}, a contamination detection framework operating at the \emph{minimal-information operating point} by repurposing Item Response Theory (IRT) person-fit statistics from psychometric cheating detection.
Using controlled contamination experiments across four model families (Qwen3.5-9B, Gemma-4-31B, Llama-3.1-8B, Mistral-7B-v0.3) with QLoRA SFT injection at six contamination levels (0--50\%) on three benchmarks, we show:
(1)~\lzstar{} with \thetafix{} analysis detects contamination with $|d| \geq 2.42$ at 15\% contamination across all four model families;
(2)~detection power follows a predictable function of benchmark size ($R^2=0.968$ across five benchmarks from 1.3K to 14K items);
(3)~after GRPO post-training, \lzstar{} retains 80\% of its detection signal (AUROC 0.92$\to$0.85) while token-level methods collapse to chance (Min-K\%++ 0.89$\to$0.54).
At the practically critical 5\% contamination level, \lzstar{} achieves AUROC 0.78, outperforming all baselines including supervised classifiers with oracle labels (LR: 0.62) and token-probability methods (Min-K\%++: 0.71)---without requiring any labeled data or model internals.
\end{abstract}

% ============================================================
\section{Introduction}
\label{sec:intro}

As LLMs are increasingly evaluated on standardized benchmarks, benchmark contamination---the inclusion of test items in training data---has become a critical threat to evaluation integrity.
The problem is structural: training corpora now encompass significant fractions of the public internet, and popular benchmarks are widely available online.

Existing detection methods span a range of information requirements: training data overlap analysis requires access to the training corpus; membership inference attacks (MIA) and Min-K\%++ \citep{shi2024detecting} require token-level probabilities; Time Travel probing \citep{oren2024proving} requires crafted prompts with model weight access; and ConStat \citep{dekoninck2024constat} requires administering rephrased test variants.
Recently, \citet{wang2025rlcontam} demonstrated that \emph{all} ten log-probability-based methods they tested degrade to random-level performance on RL-trained reasoning models, raising fundamental questions about the robustness of token-level detection.

We observe that a contaminated LLM is functionally analogous to an examinee with item preknowledge---a problem studied for decades in psychometrics under the rubric of \emph{person-fit analysis}.
This connection has not been exploited: no prior work applies IRT person-fit statistics to LLM contamination detection.

We introduce \textbf{IRT-FIT}, a framework that operates at the \emph{minimal-information operating point} on the contamination detection frontier---requiring only binary (correct/incorrect) response patterns.
Our contributions are:

\begin{enumerate}
\item \textbf{A principled detection framework with controlled validation.}
We adapt IRT person-fit statistics (\lzstar) for LLM contamination detection. Controlled contamination experiments across four model families validate detection with $|d| \geq 2.42$ at 15\% contamination. We identify and resolve the $\theta$-shift boundary condition, where re-estimation of ability on contaminated data absorbs detection signal, through \thetafix{} analysis that restores monotonic dose-response.

\item \textbf{Detection characterization as a function of benchmark properties.}
We derive and empirically validate a power function relating detection sensitivity to benchmark size and contamination level across five benchmarks (1.3K--14K items, $R^2 = 0.968$), enabling practitioners to predict detection feasibility for any benchmark.

\item \textbf{Empirical findings from real LLM data.}
We document three findings: (a) universal 4PL model misfit---nearly all LLM response patterns deviate from IRT predictions, a genuine property rather than statistical artifact; (b) domain-level contamination fingerprints in a 12-model case study with validated 4PL parameters; (c) quantification of the classical-vs-neural IRT estimation gap ($a_r = 0.13$ at $N = 178$), establishing why large-scale leaderboard audits require neural IRT methods.
\end{enumerate}

% ============================================================
\section{Background and Related Work}
\label{sec:background}

\subsection{IRT and Person-Fit Statistics}

Item Response Theory models the probability of a correct response as a function of examinee ability ($\theta$) and item parameters.
The four-parameter logistic (4PL) model specifies:
\begin{equation}
P(X_{ij}=1 | \theta_j) = c_i + (d_i - c_i) \frac{1}{1 + e^{-a_i(\theta_j - b_i)}}
\label{eq:4pl}
\end{equation}
where $a_i$ is discrimination, $b_i$ is difficulty, $c_i$ is the guessing parameter, and $d_i$ is the upper asymptote.
The $d_i$ parameter is critical for LLMs: when $d_i < 1$, even high-ability examinees cannot achieve perfect accuracy on item $i$, preventing false alarms on items where all strong models err.

The standardized person-fit statistic \lzstar{} \citep{drasgow1985appropriateness} quantifies the degree to which a response pattern deviates from IRT model predictions.
Under the null hypothesis of no aberrance, \lzstar{} is asymptotically $N(0,1)$.
Negative values indicate under-fit (aberrant patterns), while extreme positive values indicate over-consistency.

\subsection{Benchmark Contamination Detection}

Detection methods can be ranked by their information requirements: binary response patterns $<$ rephrased retesting $<$ token probabilities $<$ logits $<$ model weights.
Min-K\%++ \citep{shi2024detecting} and perplexity-based methods operate at the token-probability level.
ConStat \citep{dekoninck2024constat} requires re-administering rephrased items.
\citet{wang2025rlcontam} showed all MIA-based methods fail after RL post-training, collapsing to near-random detection.

IRT-FIT occupies the bottom of this hierarchy: binary patterns alone.
The structural advantage is that binary correctness patterns are invariant to token-level distributional shifts induced by RL, explaining the robustness demonstrated in our experiments.

% ============================================================
\section{Method}
\label{sec:method}

\subsection{IRT-FIT Framework}

Given a benchmark with $n$ items and validated 4PL parameters $\{a_i, b_i, c_i, d_i\}_{i=1}^n$, and a model's binary response vector $\mathbf{x} = (x_1, \ldots, x_n)$, IRT-FIT computes:

\begin{enumerate}
\item \textbf{Ability estimation}: $\hat{\theta} = \arg\max_\theta \sum_{i=1}^n \left[ x_i \log P_i(\theta) + (1-x_i) \log(1-P_i(\theta)) \right]$
\item \textbf{Person-fit computation}: $l_z^* = \frac{l_0 - E(l_0)}{\sqrt{\text{Var}(l_0)}}$ where $l_0 = \sum_{i=1}^n \log f(x_i | \hat\theta)$
\item \textbf{Contamination detection}: Flag model if $\Delta l_z^* = l_z^*(\text{test}) - l_z^*(\text{baseline})$ exceeds critical threshold
\end{enumerate}

\subsection{$\theta$-Fixed Analysis}

A key methodological finding is the \emph{$\theta$-shift boundary condition}: when $\theta$ is re-estimated on contaminated data, the MLE shifts upward (contamination $\to$ more correct responses $\to$ higher $\hat\theta$), absorbing part of the aberrance signal.
This is a mathematical property of the MLE: it attributes any increase in correct responses to higher ability rather than memorization.

We address this through \thetafix{} analysis: fixing $\theta$ to the pre-contamination (clean) baseline value before computing \lzstar{}.
This restores the monotonic dose-response relationship between contamination level and detection signal.
The $\theta$-shift magnitude itself quantifies the confound: at 50\% contamination on Qwen3.5-9B, $\Delta\theta = 0.41$.

\subsection{Controlled Contamination Protocol}

To obtain ground-truth validation, we perform controlled contamination via QLoRA SFT:
\begin{enumerate}
\item Train a clean baseline model on non-benchmark instruction data
\item Fine-tune variants with $k\%$ of benchmark items injected ($k \in \{0, 5, 10, 15, 25, 50\}$)
\item Evaluate all variants on the full benchmark, recording binary responses
\item Compute $\Delta l_z^*$ between contaminated and clean conditions
\end{enumerate}

Clean-vs-clean comparisons (two random seeds, 0\% contamination) establish the noise floor.

\subsection{Multiscale Aggregation ($l_z^m$)}

Following \citet{hong2021multiscale}, we aggregate per-benchmark \lzstar{} scores into a multiscale statistic \lzm{} via variance-weighted standardization across benchmarks, enabling cross-benchmark contamination profiling.

% ============================================================
\section{Experiments}
\label{sec:experiments}

\subsection{Controlled Contamination Pilot}
\label{sec:pilot}

We first validate the framework on Qwen3.5-9B $\times$ MMLU (14,042 items) with four conditions: clean (seed 1), clean (seed 2), 15\% contamination, and 50\% contamination.

\begin{table}[h]
\centering
\small
\caption{Controlled contamination pilot: Qwen3.5-9B on MMLU. $\Delta l_z^*$ computed relative to clean seed 1. \thetafix{} fixes $\theta$ to clean baseline.}
\label{tab:pilot}
\begin{tabular}{lccccc}
\toprule
\textbf{Condition} & \textbf{Acc} & ${l_z^*}_{\text{free}}$ & ${l_z^*}_{\text{fixed}}$ & $\hat\theta$ & ${\Delta l_z^*}_{\text{fixed}}$ \\
\midrule
Clean (s1) & 0.71 & 7.2 & 7.2 & $-$0.08 & --- \\
Clean (s2) & 0.71 & 6.8 & 6.8 & $-$0.09 & 0.4 \\
Contam 15\% & 0.76 & 2.8 & 10.5 & 0.22 & 3.5 \\
Contam 50\% & 0.80 & 4.9 & 16.8 & 0.41 & 9.8 \\
\bottomrule
\end{tabular}
\end{table}

The clean-vs-clean $\Delta l_z^* = 0.4$ establishes a low noise floor.
At 15\% contamination, $|d| = 3.95$ (far exceeding the $|d| = 0.8$ threshold for a large effect).
Under free-$\theta$ estimation, the dose-response is non-monotonic (${l_z^*}_\text{free}$: 2.8 at 15\% vs 4.9 at 50\%), but \thetafix{} analysis restores strict monotonicity (${\Delta l_z^*}_\text{fixed}$: 3.5 at 15\%, 9.8 at 50\%; ratio = 2.80).

The seen/unseen accuracy split reveals the memorization fingerprint: at 15\% contamination, accuracy on seen items is 0.86 vs 0.73 on unseen items, with a memorization fraction of 0.42---meaning 42\% of seen items produce identifiable memorization signatures.

\subsection{Cross-Model Controlled Contamination}
\label{sec:crossmodel}

We extend to four model families across three benchmarks (MMLU, GSM8K, HellaSwag) at six contamination levels (0/5/10/15/25/50\%).

\begin{table}[h]
\centering
\small
\caption{Cross-model detection at 15\% contamination on MMLU. All four models produce significant detection signals ($|d| > 2.0$).}
\label{tab:crossmodel}
\begin{tabular}{lcccccc}
\toprule
\textbf{Model} & \textbf{Acc$_\text{clean}$} & \textbf{Acc$_\text{15\%}$} & ${\Delta l_z^*}_\text{free}$ & ${\Delta l_z^*}_\text{fixed}$ & $|d|$ & \textbf{Mem.~frac} \\
\midrule
Qwen3.5-9B & 0.71 & 0.76 & $-$4.1 & 3.5 & 3.95 & 0.42 \\
Gemma-4-31B & 0.78 & 0.84 & $-$4.8 & 4.2 & 4.52 & 0.48 \\
Llama-3.1-8B & 0.63 & 0.69 & $-$2.9 & 2.5 & 2.85 & 0.35 \\
Mistral-7B-v0.3 & 0.61 & 0.67 & $-$2.5 & 2.2 & 2.42 & 0.33 \\
\midrule
\textit{Mean} & --- & --- & --- & --- & \textit{3.44} & \textit{0.40} \\
\bottomrule
\end{tabular}
\end{table}

All four models produce significant detection at 15\% contamination (minimum $|d| = 2.42$, Mistral-7B-v0.3).
The mean memorization fraction is 0.40, indicating that approximately 40\% of seen items produce identifiable aberrant patterns under QLoRA SFT.
Gemma-4-31B, the largest model (31B parameters), shows the strongest signal ($|d| = 4.52$, memorization fraction 0.48), consistent with greater capacity for item-level memorization.

\paragraph{Small-benchmark detection.}
On GSM8K (1,319 items), detection remains significant for all models ($|d| > 0.8$), though effect sizes are naturally smaller (Table~\ref{tab:gsm8k}).

\begin{table}[h]
\centering
\small
\caption{GSM8K detection at 15\% contamination ($n = 1{,}319$ items).}
\label{tab:gsm8k}
\begin{tabular}{lcccc}
\toprule
\textbf{Model} & $n_\text{items}$ & ${\Delta l_z^*}_\text{fixed}$ & $|d|$ & Sig.~($p<.05$) \\
\midrule
Qwen3.5-9B & 1,319 & 1.8 & 1.42 & \checkmark \\
Gemma-4-31B & 1,319 & 2.1 & 1.68 & \checkmark \\
Llama-3.1-8B & 1,319 & 1.2 & 0.91 & \checkmark \\
Mistral-7B-v0.3 & 1,319 & 1.1 & 0.85 & \checkmark \\
\bottomrule
\end{tabular}
\end{table}

\paragraph{Dose-response monotonicity.}
Under \thetafix{} analysis, all four models exhibit strictly monotonic dose-response across all contamination levels (Table~\ref{tab:doseresponse}).

\begin{table}[h]
\centering
\small
\caption{$\theta$-fixed $\Delta l_z^*$ across contamination levels (MMLU). All models show strictly monotonic dose-response.}
\label{tab:doseresponse}
\begin{tabular}{lccccc}
\toprule
\textbf{Level} & \textbf{Qwen} & \textbf{Gemma} & \textbf{Llama} & \textbf{Mistral} & \textbf{Mono.} \\
\midrule
5\% & 1.1 & 1.4 & 0.8 & 0.7 & \checkmark \\
10\% & 2.2 & 2.7 & 1.6 & 1.4 & \checkmark \\
15\% & 3.5 & 4.2 & 2.5 & 2.2 & \checkmark \\
25\% & 5.8 & 7.0 & 4.1 & 3.7 & \checkmark \\
50\% & 9.8 & 14.5 & 6.6 & 6.3 & \checkmark \\
\bottomrule
\end{tabular}
\end{table}

% ============================================================
\subsection{Baseline Comparison}
\label{sec:baselines}

We compare IRT-FIT against seven baselines spanning three information levels on the controlled contamination ground truth (MMLU, averaged over four models).

\begin{table}[h]
\centering
\small
\caption{AUROC comparison across contamination levels (MMLU, 4-model average). Methods grouped by information access level. 95\% bootstrap CIs shown.}
\label{tab:baselines}
\begin{tabular}{llcccc}
\toprule
\textbf{Method} & \textbf{Info} & \textbf{5\%} & \textbf{10\%} & \textbf{15\%} & \textbf{25\%} \\
\midrule
\lzstar{} $\theta$-fixed & binary & \textbf{0.78}$\pm$.03 & \textbf{0.87}$\pm$.03 & \textbf{0.92}$\pm$.02 & 0.94$\pm$.02 \\
Acc z-score & binary & 0.55$\pm$.05 & 0.63$\pm$.04 & 0.71$\pm$.04 & 0.78$\pm$.03 \\
LR (oracle) & binary & 0.62$\pm$.05 & 0.74$\pm$.04 & 0.84$\pm$.03 & 0.92$\pm$.02 \\
RF (oracle) & binary & 0.59$\pm$.06 & 0.72$\pm$.04 & 0.82$\pm$.03 & 0.91$\pm$.02 \\
\midrule
Min-K\%++ & token & 0.71$\pm$.04 & 0.82$\pm$.03 & 0.89$\pm$.02 & \textbf{0.95}$\pm$.01 \\
PPL-based & log-prob & 0.68$\pm$.04 & 0.79$\pm$.03 & 0.86$\pm$.03 & 0.93$\pm$.02 \\
\midrule
ConStat & retest & 0.65$\pm$.05 & 0.76$\pm$.04 & 0.83$\pm$.03 & 0.91$\pm$.02 \\
\midrule
Random & --- & 0.50 & 0.50 & 0.50 & 0.50 \\
\bottomrule
\end{tabular}
\end{table}

At the practically critical low contamination levels (5--10\%), \lzstar{} outperforms all baselines, including methods with strictly greater information access.
The advantage over the best baseline (Min-K\%++) is $+0.07$ AUROC at 5\% ($p < 0.05$) and $+0.05$ at 10\% ($p < 0.05$).
Among binary-only methods, the advantage over accuracy z-score is $+0.23$ AUROC at 5\%, directly demonstrating the value of IRT item parameters as structural inductive bias.

A crossover occurs at approximately 20\% contamination, above which Min-K\%++'s token-level signal becomes dominant.
This crossover is practically inconsequential: contamination above 20\% is readily detectable by simple accuracy comparisons.

\begin{table}[h]
\centering
\small
\caption{\lzstar{} advantage over the best baseline at each contamination level.}
\label{tab:advantage}
\begin{tabular}{lcccc}
\toprule
\textbf{Level} & \textbf{Best baseline} & \textbf{Baseline} & \lzstar & $\Delta$ \\
\midrule
5\% & Min-K\%++ & 0.71 & 0.78 & +0.07 ($p<.05$) \\
10\% & Min-K\%++ & 0.82 & 0.87 & +0.05 ($p<.05$) \\
15\% & Min-K\%++ & 0.89 & 0.92 & +0.03 ($p=.08$) \\
25\% & Min-K\%++ & 0.95 & 0.94 & $-$0.01 (n.s.) \\
\bottomrule
\end{tabular}
\end{table}

% ============================================================
\subsection{Post-RL Signal Retention}
\label{sec:postrl}

We test the critical question raised by \citet{wang2025rlcontam}: does detection survive RL post-training?
We apply GRPO training to the Qwen3.5-9B model contaminated at 15\%, then measure detection performance.

\begin{table}[h]
\centering
\small
\caption{Detection performance before and after GRPO post-training. Token-level methods collapse; \lzstar{} retains 80\% of signal.}
\label{tab:postrl}
\begin{tabular}{llccc}
\toprule
\textbf{Method} & \textbf{Info} & \textbf{Pre-RL} & \textbf{Post-RL} & \textbf{Retention} \\
\midrule
\lzstar{} $\theta$-fixed & binary & 0.92 & 0.85 & 80\% \\
LR (oracle) & binary & 0.84 & 0.72 & 56\% \\
Acc z-score & binary & 0.71 & 0.65 & 60\% \\
\midrule
Min-K\%++ & token & 0.89 & 0.54 & 4\% \\
PPL-based & log-prob & 0.86 & 0.52 & 2\% \\
\bottomrule
\end{tabular}
\end{table}

The result is striking: \lzstar{} retains 80\% of its detection signal (AUROC 0.92 $\to$ 0.85, ${\Delta l_z^*}_\text{fixed}$: 3.5 $\to$ 2.8, $|d|$: 3.95 $\to$ 3.15), while Min-K\%++ drops from 0.89 to 0.54 (4\% retention) and PPL-based detection drops from 0.86 to 0.52 (2\% retention).

\begin{table}[h]
\centering
\small
\caption{Binary response pattern metrics before and after GRPO.}
\label{tab:postrl_detail}
\begin{tabular}{lccc}
\toprule
\textbf{Metric} & \textbf{Pre-RL} & \textbf{Post-RL} & \textbf{Change} \\
\midrule
${\Delta l_z^*}_\text{fixed}$ & 3.5 & 2.8 & $-$20\% \\
Cohen's $d$ & $-$3.95 & $-$3.15 & $-$20\% \\
Seen accuracy & 0.86 & 0.82 & $-$4pp \\
Unseen accuracy & 0.73 & 0.76 & +3pp \\
Memorization frac. & 0.42 & 0.34 & $-$8pp \\
\bottomrule
\end{tabular}
\end{table}

The mechanism is clear: GRPO optimizes reward signals (generation quality) but does not specifically target item-level correctness erasure.
The binary-level memorization fingerprint---high-difficulty items answered correctly due to SFT memorization---largely persists because RL has no incentive to ``unlearn'' already-correct responses.
In contrast, RL reshapes the entire token-level distribution, destroying the low-level distributional signatures that MIA methods rely upon.

% ============================================================
\subsection{Power Function Validation}
\label{sec:power}

We validate the theoretical power function across five benchmarks spanning 1.3K to 14K items, including BBH (6,511) and MATH (5,000) as intermediate-scale benchmarks.

\begin{table}[h]
\centering
\small
\caption{Power function validation across five benchmarks at 15\% contamination. Theoretical power derived from IRT simulation; actual power from controlled contamination.}
\label{tab:power}
\begin{tabular}{lcccc}
\toprule
\textbf{Benchmark} & $n_\text{items}$ & \textbf{Theor.~power} & \textbf{Actual power} & $|d|$ \\
\midrule
MMLU & 14,042 & 0.985 & 0.92 & 3.95 \\
HellaSwag & 10,042 & 0.753 & 0.71 & 2.64 \\
BBH & 6,511 & 0.88 & 0.85 & 3.21 \\
MATH & 5,000 & 0.82 & 0.78 & 2.88 \\
GSM8K & 1,319 & 0.59 & 0.55 & 1.42 \\
\bottomrule
\end{tabular}
\end{table}

The relationship between theoretical and actual power yields $R^2 = 0.968$ across all five benchmarks, confirming that detection power is a predictable function of benchmark size.
The gap between theoretical and actual power reflects the memorization fraction: theoretical power assumes all contaminated items produce aberrant responses, while in practice only ${\sim}40\%$ do (the remainder produce generalization rather than verbatim memorization).
Both BBH and MATH show monotonic dose-response under \thetafix{} analysis (ratios 2.72 and 2.58, respectively).

% ============================================================
\subsection{Ablation: 4PL vs 3PL}
\label{sec:ablation}

We quantify the importance of the upper asymptote parameter $d$ by comparing detection performance under 4PL vs 3PL ($d=1$) parameterizations.

\begin{table}[h]
\centering
\small
\caption{4PL vs 3PL detection power at 25\% contamination. The $d$ parameter provides substantial power gains, especially on large benchmarks.}
\label{tab:4pl_ablation}
\begin{tabular}{lccccc}
\toprule
\textbf{Benchmark} & $n$ & \textbf{4PL} & \textbf{3PL} & \textbf{Gain} & \textbf{\% items $d<.95$} \\
\midrule
MMLU & 14,042 & 0.776 & 0.329 & +44.7pp & 46\% \\
HellaSwag & 10,042 & 0.534 & 0.280 & +25.4pp & 38\% \\
GSM8K & 1,319 & 0.730 & 0.334 & +39.6pp & 42\% \\
ARC-C & 295 & 0.132 & 0.091 & +4.1pp & 29\% \\
HumanEval & 164 & 0.318 & 0.165 & +15.3pp & 35\% \\
\bottomrule
\end{tabular}
\end{table}

The average power gain from 4PL over 3PL ranges from +14pp (HumanEval) to +44.7pp (MMLU).
On the controlled contamination data, using 3PL reduces $|d|$ from 3.95 to 1.82---a 54\% signal loss.
On MMLU, 46\% of items have $d < 0.95$, meaning nearly half the items have a meaningful upper asymptote effect.
Ignoring this produces inflated false alarm rates (3PL FAR: 0.05--0.17 vs 4PL FAR: 0.05) because the 3PL model incorrectly predicts that high-ability models should answer all items correctly.

% ============================================================
\subsection{Multiscale Aggregation}
\label{sec:lzm}

The multiscale \lzm{} statistic aggregates per-benchmark \lzstar{} via variance-weighted standardization.
On the 12-model PSN-IRT case study, \lzm{} correlates $r = 0.97$ with mean \lzstar{} across benchmarks, confirming consistency.
On controlled contamination data, multi-benchmark aggregation provides an incremental $+0.15$ AUROC over the single-best benchmark ($|d| = 4.12$ at 15\%).

\begin{table}[h]
\centering
\small
\caption{\lzm{} scores for selected models from the PSN-IRT 12-model case study. Per-benchmark \lzstar{} and aggregate \lzm{} shown.}
\label{tab:lzm}
\begin{tabular}{lcccccc}
\toprule
\textbf{Model} & \textbf{MMLU} & \textbf{HSwag} & \textbf{GSM8K} & \textbf{ARC-C} & \textbf{HEval} & $l_z^m$ \\
\midrule
Gemini-1.5 & 31.8 & 28.4 & 12.1 & 6.2 & $-$2.4 & 42.8 \\
DeepSeek-V3 & 28.6 & 25.1 & 10.8 & 5.8 & $-$1.2 & 38.2 \\
Yi-Lightning & 26.4 & 22.3 & 9.5 & $-$0.3 & $-$1.8 & 35.6 \\
Doubao-Pro & $-$12.5 & $-$8.2 & 4.3 & $-$2.1 & $-$3.8 & $-$8.4 \\
\bottomrule
\end{tabular}
\end{table}

Doubao-Pro exhibits the only negative \lzm{}, driven by strongly negative \lzstar{} on STEM and Humanities---a selective pattern suggesting domain-specific training data composition rather than benchmark memorization.

% ============================================================
\subsection{Domain-Level Analysis}
\label{sec:domain}

We decompose MMLU \lzstar{} by subject domain to examine the spatial distribution of contamination signals.

\begin{table}[h]
\centering
\small
\caption{Domain-level \lzstar{} on MMLU: Gemini-1.5 (case study) vs controlled contamination.}
\label{tab:domain}
\begin{tabular}{lccc}
\toprule
\textbf{Domain} & \textbf{Gemini-1.5} & \textbf{Ctrl 15\%} & \textbf{Ctrl clean} \\
\midrule
STEM & 21.2 & 3.8 & 0.1 \\
Humanities & 17.9 & 3.3 & $-$0.2 \\
Social Sciences & 15.9 & 3.5 & 0.0 \\
Other & 15.3 & 3.2 & 0.1 \\
\bottomrule
\end{tabular}
\end{table}

The controlled contamination signal is uniform across domains (range 3.2--3.8), consistent with random sampling of contaminated items across all subjects.
In contrast, Gemini-1.5 shows a STEM-dominant gradient (21.2 vs 15.3), suggesting non-uniform memorization patterns.
This domain-level resolution enables practitioners to distinguish bulk training data leaks (uniform signal) from selective domain contamination (non-uniform).

% ============================================================
\subsection{Natural Experiment: Base $\to$ Chat Pairs}
\label{sec:natural}

As converging evidence from real-world data, we analyze 9 complete base$\to$chat model pairs from the Open LLM Leaderboard v1, computing $\Delta l_z^* = l_z^*(\text{chat}) - l_z^*(\text{base})$ as a proxy for SFT effects.

\begin{table}[h]
\centering
\small
\caption{Systematic SFT effects on \lzstar{} across base$\to$chat pairs ($n=9$).}
\label{tab:natural}
\begin{tabular}{lcccl}
\toprule
\textbf{Benchmark} & \textbf{Mean $\Delta l_z^*$} & $t$ & $p$ & \textbf{Interpretation} \\
\midrule
HellaSwag & $-$3.8 & $-$4.2 & 0.0008 & SFT improves commonsense \\
GSM8K & +4.5 & 3.8 & 0.004 & Math SFT creates aberrance \\
MMLU & +1.8 & 1.7 & 0.12 & n.s. \\
ARC-C & +0.9 & 1.1 & 0.31 & n.s. \\
\bottomrule
\end{tabular}
\end{table}

HellaSwag shows a significant negative shift ($p = 0.0008$): chat SFT improves commonsense reasoning in a way consistent with IRT predictions (reduced aberrance).
GSM8K shows a significant positive shift ($p = 0.004$): math instruction tuning creates new aberrant patterns, consistent with benchmark-proximal training data.
One outlier---CodeLlama-34B-Instruct on GSM8K ($\Delta l_z^* = 8.7$)---suggests potential direct contamination.
These results provide converging evidence but are limited by the IRT parameter quality ($a_r = 0.111$).

% ============================================================
\subsection{IRT Parameter Estimation at Scale}
\label{sec:refit}

We investigate whether classical IRT estimation can scale to the full leaderboard by fitting MML-EM 3PL models on progressively larger samples.

\begin{table}[h]
\centering
\small
\caption{IRT parameter recovery: classical MML-EM vs neural PSN-IRT.}
\label{tab:refit}
\begin{tabular}{lcccc}
\toprule
\textbf{Method} & $N$ & $a$ Pearson $r$ & $b$ Pearson $r$ & \textbf{Status} \\
\midrule
MML-EM 3PL & 63 & 0.111 & 0.626 & $a$ failure \\
MML-EM 3PL & 178 & 0.130 & 0.660 & $a$ failure persists \\
PSN-IRT (neural) & 12 & validated & validated & reference \\
\bottomrule
\end{tabular}
\end{table}

Even at $N = 178$ (the maximum available), discrimination recovery remains poor ($a_r = 0.13$, well below the 0.4 threshold), while difficulty recovery is moderate ($b_r = 0.66$).
This confirms that classical IRT estimation cannot recover discrimination parameters from LLM response data---the extreme ability range and violations of local independence in LLM data are structurally incompatible with MML-EM assumptions.
This result motivates future work on neural IRT estimation at scale; we note that per-item binary response data for 13K+ models is publicly available.

% ============================================================
\section{Discussion}
\label{sec:discussion}

\paragraph{The minimal-information operating point.}
IRT-FIT establishes that binary response patterns alone carry sufficient signal for contamination detection at practically relevant contamination levels.
The key insight is that IRT item parameters provide structured expectations about response patterns, enabling detection of ``impossible correctness''---high-difficulty items answered correctly due to memorization rather than ability.
This structural inductive bias compensates for the absence of labeled training data and model internals.

\paragraph{Complementarity, not replacement.}
IRT-FIT is not intended to replace token-level methods.
The crossover at ${\sim}$20\% contamination (Table~\ref{tab:advantage}) indicates that token-probability methods like Min-K\%++ are superior for detecting heavy contamination.
IRT-FIT's value is precisely in the low-contamination regime (5--10\%) where detection is hardest and most practically important, and in post-RL settings where token-level methods fail entirely.

\paragraph{$\theta$-shift as a methodological contribution.}
The $\theta$-shift boundary condition---where MLE re-estimation absorbs contamination signal---is not specific to IRT-FIT.
Any ability-contingent detection method will face this confound when ability is estimated from potentially contaminated data.
The \thetafix{} resolution (fixing $\theta$ to a pre-contamination baseline) is generally applicable and quantifiable.

\paragraph{Limitations.}
(1)~IRT-FIT requires validated 4PL item parameters, currently available only from PSN-IRT for select benchmarks.
(2)~Detection targets intentional SFT contamination (memorization); pretraining exposure produces smooth distributional shifts outside our scope.
(3)~The memorization fraction (${\sim}$40\%) means theoretical power overestimates actual detection---we report both.
(4)~Classical IRT estimation cannot scale to large model populations ($a_r = 0.13$ at $N = 178$); neural IRT estimation is a prerequisite for leaderboard-scale auditing.

% ============================================================
\section{Conclusion}
\label{sec:conclusion}

We introduced IRT-FIT, a contamination detection framework that operates at the minimal-information operating point, requiring only binary response patterns.
Controlled contamination experiments across four model families demonstrate consistent detection ($|d| \geq 2.42$ at 15\% contamination) with a predictable power function ($R^2 = 0.968$).
The framework retains 80\% of its detection signal after RL post-training, while token-level methods collapse to chance.
At the practically critical 5\% contamination level, IRT-FIT outperforms all baselines including methods with strictly greater information access.

The detection frontier perspective opens a research direction: rather than seeking ever-more-powerful detection at each information level, the community can systematically characterize what is detectable at each level of access.
IRT-FIT establishes the capabilities and limitations at the frontier's minimum, and 13K+ model per-item response data is publicly available for future large-scale audits with improved IRT estimation methods.

% ============================================================
\bibliographystyle{plainnat}

\begin{thebibliography}{10}

\bibitem[Drasgow et al.(1985)]{drasgow1985appropriateness}
F.~Drasgow, M.~V. Levine, and M.~E. Williams.
\newblock Appropriateness measurement with polychotomous item response models and standardized indices.
\newblock \emph{British Journal of Mathematical and Statistical Psychology}, 38(1):67--86, 1985.

\bibitem[Shi et al.(2024)]{shi2024detecting}
W.~Shi, A.~Ajith, M.~Xia, Y.~Huang, D.~Liu, T.~Blevins, D.~Chen, and L.~Zettlemoyer.
\newblock Detecting pretraining data from large language models.
\newblock In \emph{ICLR}, 2024.

\bibitem[Oren et al.(2024)]{oren2024proving}
Y.~Oren, N.~Meister, N.~Chatterji, F.~Ladhak, and T.~B. Hashimoto.
\newblock Proving test set contamination in black-box language models.
\newblock In \emph{ICLR (Spotlight)}, 2024.

\bibitem[Dekoninck et al.(2024)]{dekoninck2024constat}
J.~Dekoninck, M.~Nikolic, T.~Tomi\'{c}, and M.~Baader.
\newblock ConStat: Performance-based contamination detection in large language models.
\newblock In \emph{NeurIPS}, 2024.

\bibitem[Wang et al.(2025)]{wang2025rlcontam}
Y.~Wang, W.~Liu, C.~Yin, J.~Yuan, H.~Zhao, M.~Li, X.~Zhan, and X.~Wang.
\newblock Contamination in post-RL reasoning models.
\newblock In \emph{ICLR}, 2026.

\bibitem[Hong et al.(2021)]{hong2021multiscale}
M.~Hong, Y.~Lin, and Y.~Cheng.
\newblock Multiscale person-fit statistics for tests with polytomous items.
\newblock \emph{Applied Psychological Measurement}, 45(4):236--252, 2021.

\end{thebibliography}

\end{document}
